\documentclass[12pt]{article}
\usepackage[utf8]{inputenc}
\usepackage{geometry}
\geometry{letterpaper, margin=0.25in}
\usepackage{graphicx} 
\usepackage{multicol}
\usepackage{parskip}
\usepackage{booktabs}
\usepackage{array} 
\usepackage{paralist} 
\usepackage{verbatim}
\usepackage{sectsty}
\usepackage[shortlabels]{enumitem}

\usepackage{amsmath}
\usepackage{amssymb}
\usepackage{mathtools}
\usepackage{empheq}
\usepackage{xcolor}
\usepackage{bbm}
\usepackage{tikz}
\usepackage{pgfplots}
\usepackage{tikz-cd}
\usepackage{circuitikz}
\usepackage{multirow}
\pgfplotsset{compat=1.18}
\usetikzlibrary{intersections, decorations.markings}
\tikzset{
    marking along/.style n args={2}{
        decoration={
                markings, 
                mark=at position #1 with {\arrow{#2}}
        },
        postaction={decorate}
        },
    marking along/.default={0.5}{>}
    wavy/.style={
        decorate,decoration={coil,aspect=0}
     },
    two marks/.style n args={1}{
        decoration={
            markings,
            mark=at position 0.25 with {\arrow{#1}},
            mark=at position 0.75 with {\arrow{#1}}
         },
         postaction={decorate}
    }
    two marks/.default={>}
}

\colorlet{mygreen}{green!50!teal}
\colorlet{darkblue}{blue!60!black}

\newcommand{\ans}[1]{\boxed{\text{#1}}}
\newcommand{\vecs}[1]{\langle #1\rangle}
\renewcommand{\hat}[1]{\widehat{#1}}

\renewcommand{\P}{\mathbb{P}}
\newcommand{\R}{\mathbb{R}}
\newcommand{\E}{\mathbb{E}}
\newcommand{\Z}{\mathbb{Z}}
\newcommand{\N}{\mathbb{N}}
\newcommand{\Q}{\mathbb{Q}}
\newcommand{\C}{\mathbb{C}}

\newcommand{\ind}{\mathbbm{1}}
\newcommand{\qed}{\quad \blacksquare}

\newcommand{\brak}[1]{\left\langle #1 \right\rangle}
\newcommand{\bra}[1]{\left\langle #1 \right\vert}
\newcommand{\ket}[1]{\left\vert #1 \right\rangle}

\newcommand{\abs}[1]{\left\vert #1 \right\vert}
\newcommand{\mfX}{\mathfrak{X}}
\newcommand{\ep}{\varepsilon}

\newcommand{\Ec}{\mathcal{E}}
\newcommand{\Nc}{\mathcal{N}}
\newcommand{\A}{\mathcal{A}}
\newcommand{\F}{\mathcal{F}}
\newcommand{\Cc}{\mathcal{C}}
\newcommand{\B}{\mathcal{B}}
\newcommand{\M}{\mathcal{M}}
\newcommand{\X}{\chi}
\renewcommand{\L}{\mathcal{L}}

\newcommand{\sub}{\subseteq}
\newcommand{\st}{\text{ s.t. }}
\newcommand{\card}{\text{card }}
\renewcommand{\div}{\vspace*{10pt}\hrule\vspace*{10pt}}
\newcommand{\surj}{\twoheadrightarrow}
\newcommand{\inj}{\hookrightarrow}
\newcommand{\biject}{\hookrightarrow \hspace{-8pt} \rightarrow}
\renewcommand{\bar}[1]{\overline{#1}}
\newcommand{\overcirc}[1]{\overset{\circ}{#1}}
\newcommand{\diam}{\text{diam }}
\newcommand{\iid}{\overset{	ext{iid}}{\sim}}

\renewcommand{\Re}{\text{Re}\,}
\renewcommand{\Im}{\text{Im}\,}
\newcommand{\Var}{\text{Var}\,}
\newcommand{\Cov}{\text{Cov}\,}

\DeclareMathOperator*{\argmax}{\arg\max}
\DeclareMathOperator*{\argmin}{\arg\min}

\newcommand{\sign}{\text{sign}\,}

\newcommand*{\tbf}[1]{\ifmmode\mathbf{#1}\else\textbf{#1}\fi}

\usepackage{tcolorbox}
\tcbuselibrary{breakable, skins}
\tcbset{enhanced}
\newenvironment*{tbox}[2][gray]{
    \begin{tcolorbox}[
        parbox=false,
        colback=#1!5!white,
        colframe=#1!75!black,
        breakable,
        title={#2}
    ]}
    {\end{tcolorbox}}

\newenvironment*{exercise}[1][red]{
    \begin{tcolorbox}[
        parbox=false,
        colback=#1!5!white,
        colframe=#1!75!black,
        breakable
    ]}
    {\end{tcolorbox}}

\newenvironment*{proof}[1][blue]{
\begin{tcolorbox}[
    parbox=false,
    colback=#1!5!white,
    colframe=#1!75!black,
    breakable
]}
{\end{tcolorbox}}
\newenvironment*{proposition}[1][gray]{
    \begin{tcolorbox}[
        parbox=false,
        colback=#1!5!white,
        colframe=#1!75!black,
        breakable
    ]}
    {\end{tcolorbox}}


\begin{document}
\begin{multicols}{2}
    \begin{flushleft}
        \parbox{0.5\textwidth}{\Large ENGN 1580 - Homework 4}
    \end{flushleft}
    \begin{flushright}
        \tbf{Milan Capoor}\\
        Spring 2026
    \end{flushright}
\end{multicols}
\div

\tbf{1. Counting and Sample Space:} Consider a binary code with 5 bits (0 or 1) in each code
word. An example of a code word is 01010. In each code word, a bit is zero with probability
0.8, independent of any other bit.
\begin{enumerate}
\item What is the probability of the code word 00111?

\color{darkblue}
    \[(0.8)^2 (0.2)^3\]
\color{black}

\item What is the probability that a code word contains exactly three ones?

\color{darkblue}
\[\binom{5}{3} (0.8)^2 (0.2)^3\]

\color{black}
\end{enumerate}

\tbf{2. Simple Probability, Simple Application:} A source wishes to transmit radio packets to a receiver over a radio link. The receiver uses error detection to identify packets that have been
corrupted by radio noise. When a packet is sent error free the receiver sends an acknowledge-
ment(ACK) back to the source. When the receiver gets a packet with errors, it sends back
a negative-acknowledgement(NACK). Each time the source receives a NACK, the packet
is re-transmitted. We assume that each packet transmission is independently corrupted by
errors with probability $q$.
\begin{enumerate}

\item Find the PMF of $X$, the number of times that a packet is transmitted by the source.

\color{darkblue}
    \[f_X(x) = \begin{cases}
        q^{x-1} \cdot (1-q) & x > 0\\ 
        0 & x \leq 0
    \end{cases}\]
\color{black}   

\item Suppose each packet takes 1 millisecond to transmit and that the source waits an additional millisecond to receive the acknowledgement message (ACK or NACK) before
transmitting. Let $T$ equal the time required until the packet is successfully received.
What is the relationship between $T$ and $X$? What is the PMF of $T$?


\end{enumerate}

\tbf{3. Gaussians:} W is a Gaussian random variable with expected value $\mu = 0$ and variance
$\sigma^2 = 16$. Given the event $C = {W > 0}$,
\begin{enumerate}

\item What is the conditional PDF $f_{W |C} (w)$?

\color{darkblue}
    \[W \sim \Nc(0, 16)\]
\color{black}

\item What is the conditional expected value $\E[W |C]$?

\color{darkblue}
    \[\E[W | C] = \int_{-\infty}^{\infty} \frac{1}{\sqrt{8\pi}} e^{-x^2/32} \; dx = \int_0^{\infty} \frac{1}{\sqrt{8\pi}} e^{-x^2/32} \; dx  =\frac{32}{\sqrt{8 \pi}}\]
\color{black}

\item \textcolor{red}{Find the conditional variance, $\Var[W |C]$.}


\end{enumerate}

\tbf{4. Functions of random variables:}
\begin{enumerate}

\item Find the PDF $f_Y (y)$ for the random variable $Y$, where $Y = X^3$ and $X$ is a uniform
random variable in the range (0,1).

\color{darkblue}
\begin{align*}
    F_Y(y) &= \P(Y \leq y) = \P(X^3 \leq y) = \P(X \leq \sqrt[3]{y}) = F_X(\sqrt[3]{y}) = \sqrt[3]{y}\\ 
    f_Y(y) &= \frac{1}{3}y^{-2/3}
\end{align*}
\color{black}

\item Let X and Y be independent random variables with $f_X (x) = \exp(-x)u(x)$, and
$f_Y (y) = \frac{1}{2} [u(y + 1) - u(y - 1)]$ and let $Z = X + Y$, where $u(\cdot)$ denotes unit step function. What is the PDF of random variable $Z$? HINT:Use convolution!

\color{darkblue}
\begin{align*}
    f_Z(z) &= f_X \star f_Y = \int_{0}^1 \frac{1}{2}\exp(-x) = -\frac{1}{2e} + \frac{1}{2} 
\end{align*}
\color{black}
\end{enumerate}

\tbf{5. Gaussian + Gaussian = Gaussian:} PROVE that if $X$ and $Y$ are two zero mean unit variance
independent Gaussian random variables that the density of $Z = X + Y$ is a zero mean
Gaussian random variable with variance 2.

\color{darkblue}
    \begin{align*}
        f_Z(z) &= \int \frac{1}{\sqrt{2\pi}} e^{-u^2/2} \cdot \frac{1}{\sqrt{2\pi}} e^{-(z - u)^2/2} \; du\\ 
        &= \frac{1}{2\pi} \int \exp\left(\frac{1}{2} \cdot (-u^2 - z^2 + 2zu - u^2 )  \right)\; du\\ 
        &= \textcolor{red}{\frac{1}{2\pi} \exp\left(-\frac{z^2}{2}\right) \int \exp\left(-u^2 - zu \right)\; du}\\ 
    \end{align*}
\color{black}

\end{document}